% !TeX root = ../../Book2.tex
% 纲要标题：### 10.6 非零 d₂ 的双复形模型
% 纲要位置：几何学计划纲要.md，第 1423 行。
% 正文状态：初稿；尚未正式审核。

% 纲要内容开始
% 
% 取第一象限双复形 C^{0,1}=C^{1,0}=C^{1,1}=C^{2,0}=k，令从(1,0)到(1,1)、从(0,1)到(1,1)、从(1,0)到(2,0)的三个微分为恒等，其余为零。逐页计算、提升代表元并直接核验总复形无上同调；符号按本卷反交换约定统一。此节是一段连续计算，不再人为分小节。
% 
% ---
% 
% 纲要内容结束

\section{非零 \texorpdfstring{\(\mathrm{d}_2\)}{d₂} 的双复形模型}
\label{g2:sec:1006}

高阶微分可以出现在只有水平、竖直两种微分的双复形中。关键在于一个横向像虽然在纵向上同调中为零，用来消去它的纵向原像还可能产生新的横向像。下面把这个过程放在四个一维向量空间中，从零页一直算到总上同调。

\ProseParagraphBreak

设 \(k\) 为任意域，双复形的非零项为
\[
	C^{0,1}=ku,\qquad C^{1,0}=kv,\qquad C^{1,1}=kw,\qquad C^{2,0}=kz.
\]
选定各自的基，规定
\begin{equation}\label{g2:eq:nonzero-d2-model}
	\mathrm d_{\mathrm h}u=w,\qquad \mathrm d_{\mathrm v}v=w,\qquad \mathrm d_{\mathrm h}v=z,
\end{equation}
其余微分均为零。三个非零箭头在这些基下都是恒等映射。把上排视为 \(q=1\)、下排视为 \(q=0\)，三列依次为 \(p=0,1,2\)，图为
\[
	\begin{tikzcd}[column sep=large,row sep=large]
		ku \arrow[r,"1"] & kw & 0 \\
		0 & kv \arrow[u,"1"] \arrow[r,"1"'] & kz.
	\end{tikzcd}
\]
任意两个可复合的微分都为零：一旦到达 \(w\) 或 \(z\)，已经没有非零的后续箭头。因此 \(\mathrm d_{\mathrm h}^2=\mathrm d_{\mathrm v}^2=0\) 且 \(\mathrm d_{\mathrm h}\mathrm d_{\mathrm v}+\mathrm d_{\mathrm v}\mathrm d_{\mathrm h}=0\)，这确实是本章约定的反交换双复形。

\ProseParagraphBreak

先把目标算出来，便于之后逐页核对。总复形只有两个非零次数：
\begin{equation}\label{g2:eq:nonzero-d2-total-complex}
	0\longrightarrow ku\oplus kv\xrightarrow{\mathrm d}kw\oplus kz\longrightarrow0,
	\qquad [\mathrm d]=\begin{pmatrix}1&1\\0&1\end{pmatrix}.
\end{equation}
左项在次数 \(1\)，右项在次数 \(2\)，因为 \(\mathrm du=w\)、\(\mathrm dv=w+z\)。该矩阵行列式为 \(1\)，所以微分是同构，\(H^n(\operatorname{Tot}C)=0\) 对所有 \(n\) 成立。谱序列也应当最后把所有类消去。

\ProseParagraphBreak

按列取递减滤过 \(F^p\operatorname{Tot}^nC=\bigoplus_{a\ge p}C^{a,n-a}\)。在两个非零次数，它分别是
\[
	F^0\operatorname{Tot}^1=ku\oplus kv,\quad F^1\operatorname{Tot}^1=kv,\quad F^2\operatorname{Tot}^1=0,
\]
\[
	F^0\operatorname{Tot}^2=F^1\operatorname{Tot}^2=kw\oplus kz,\quad F^2\operatorname{Tot}^2=kz,\quad F^3\operatorname{Tot}^2=0.
\]
这些式子也直接验证了总微分保持滤过。

\ProseParagraphBreak

零页就是原来的四项双分次向量空间，\(\mathrm d_0=\mathrm d_{\mathrm v}\)。中间列的 \(kv\xrightarrow{1}kw\) 是同构，取列上同调后两个位置都消失；其余两列各留下一项。因此
\[
	E_1^{0,1}=k[u]_1,\qquad E_1^{2,0}=k[z]_1,
\]
其他项全为零。第一页微分 \(\mathrm d_1\) 的双次数为 \((1,0)\)，两个非零位置之间没有这种箭头，所以 \(\mathrm d_1=0\)、\(E_2=E_1\)。用代表元看，\(\mathrm d_{\mathrm h}u=w\) 在中间列是纵向边界，故 \(\mathrm d_1[u]_1=0\)。

\ProseParagraphBreak

这个消失允许我们修正 \(u\)。取 \(x_1=-v\)，就有
\[
	\mathrm d_{\mathrm h}u+\mathrm d_{\mathrm v}(-v)=w-w=0.
\]
于是修正后的总上链 \(x=u-v\) 满足
\begin{equation}\label{g2:eq:nonzero-d2-corrected-representative}
	\mathrm d(u-v)=w-(w+z)=-z\in F^2\operatorname{Tot}^2C.
\end{equation}
因此它属于 \(Z_2^{0,1}\)，是第二页中从 \([u]_1\) 延续下来的类。按\cref{g2:prop:spectral-sequence-differential}，
\begin{equation}\label{g2:eq:nonzero-d2-value}
	\mathrm d_2:E_2^{0,1}\longrightarrow E_2^{2,0},\qquad [u-v]_2\longmapsto-[z]_2.
\end{equation}
两边均为一维，\(-1\) 在任意域中都是单位；特征 \(2\) 时它等于 \(1\)，映射仍然非零。因此 \(\mathrm d_2\) 是同构，第三页全部为零，此后各页也全为零。

\ProseParagraphBreak

这项微分的两个作用都能在总复形中看见。源中的 \(u\) 无法改成总闭元：要消掉 \(w\) 分量，修正系数必须是 \(-v\)，但剩下的 \(-z\) 已经没有更深的次数 \(1\) 上链可供抵消，因为 \(F^2\operatorname{Tot}^1=0\)。靶中的 \(z\) 本来是一个总闭元，却满足 \(z=\mathrm d(v-u)\)，所以它是总边界；这个边界的原像位于较浅的滤过层，直到第二页才进入当前层的关系中。一张 \(\mathrm d_2\) 同时删除了不能继续提升的源类和本来就应被识别为边界的靶类。

\ProseParagraphBreak

全部页面可归纳为
\[
	\begin{array}{c|c|c}
		\text{页面}&\text{非零位置}&\text{非零微分}\\ \hline
		E_0&(0,1),(1,0),(1,1),(2,0)&\mathrm d_0(v)=w\\
		E_1&(0,1),(2,0)&\text{无}\\
		E_2&(0,1),(2,0)&\mathrm d_2[u-v]_2=-[z]_2\\
		E_3=E_\infty&\text{无}&\text{无}
	\end{array}
\]
第一象限收敛定理给出的 \(E_\infty=\operatorname{gr}_F H^*(\operatorname{Tot}C)=0\)，与\eqref{g2:eq:nonzero-d2-total-complex}的直接计算吻合。若改按行滤过，上排 \(ku\to kw\) 和下排 \(kv\to kz\) 都已经是同构，于是另一条谱序列在第一页就为零。两种滤过改变的是计算中何时看见这些抵消，总复形及其上同调始终相同。
% 本节沿用纲要的连续叙述安排，不另造小节标题。

